% Options for packages loaded elsewhere
\PassOptionsToPackage{unicode}{hyperref}
\PassOptionsToPackage{hyphens}{url}
\PassOptionsToPackage{dvipsnames,svgnames,x11names}{xcolor}
%
\documentclass[
  11pt,
  a4paper,
]{ltjsarticle}
\usepackage{amsmath,amssymb}
\usepackage{iftex}
\ifPDFTeX
  \usepackage[T1]{fontenc}
  \usepackage[utf8]{inputenc}
  \usepackage{textcomp} % provide euro and other symbols
\else % if luatex or xetex
  \usepackage{unicode-math} % this also loads fontspec
  \defaultfontfeatures{Scale=MatchLowercase}
  \defaultfontfeatures[\rmfamily]{Ligatures=TeX,Scale=1}
\fi
\usepackage{lmodern}
\ifPDFTeX\else
  % xetex/luatex font selection
\fi
% Use upquote if available, for straight quotes in verbatim environments
\IfFileExists{upquote.sty}{\usepackage{upquote}}{}
\IfFileExists{microtype.sty}{% use microtype if available
  \usepackage[]{microtype}
  \UseMicrotypeSet[protrusion]{basicmath} % disable protrusion for tt fonts
}{}
\makeatletter
\@ifundefined{KOMAClassName}{% if non-KOMA class
  \IfFileExists{parskip.sty}{%
    \usepackage{parskip}
  }{% else
    \setlength{\parindent}{0pt}
    \setlength{\parskip}{6pt plus 2pt minus 1pt}}
}{% if KOMA class
  \KOMAoptions{parskip=half}}
\makeatother
\usepackage{xcolor}
\usepackage[margin=24mm]{geometry}
\setlength{\emergencystretch}{3em} % prevent overfull lines
\providecommand{\tightlist}{%
  \setlength{\itemsep}{0pt}\setlength{\parskip}{0pt}}
\setcounter{secnumdepth}{5}
\ifLuaTeX
  \usepackage{selnolig}  % disable illegal ligatures
\fi
\IfFileExists{bookmark.sty}{\usepackage{bookmark}}{\usepackage{hyperref}}
\IfFileExists{xurl.sty}{\usepackage{xurl}}{} % add URL line breaks if available
\urlstyle{same}
\hypersetup{
  colorlinks=true,
  linkcolor={blue},
  filecolor={Maroon},
  citecolor={Blue},
  urlcolor={blue},
  pdfcreator={LaTeX via pandoc}}

\author{}
\date{}

\begin{document}

\hypertarget{chapter-09-ux4f4eux30e9ux30f3ux30afux69cbux9020pcaux5727ux7e2e}{%
\section{Chapter 09 ---
低ランク構造・PCA・圧縮}\label{chapter-09-ux4f4eux30e9ux30f3ux30afux69cbux9020pcaux5727ux7e2e}}

この章では「高次元に見えるデータが、実は少数の方向でほとんど説明できる」という考えを扱います。

\hypertarget{rank-ux306fux81eaux7531ux5ea6ux3092ux6e2cux308b}{%
\subsection{rank
は自由度を測る}\label{rank-ux306fux81eaux7531ux5ea6ux3092ux6e2cux308b}}

行列の rank が \texttt{r} なら、その列や行に含まれる独立方向は
\texttt{r} 個だけです。

現実のデータは厳密には full rank でも、特異値が

\[
100,20,5,0.1,0.01,\dots
\]

のように急減することがあります。この場合「近似的には低ランク」と考えられます。

\hypertarget{svd-ux304cux6700ux826fux8fd1ux4f3cux3092ux4e0eux3048ux308b}{%
\subsection{SVD
が最良近似を与える}\label{svd-ux304cux6700ux826fux8fd1ux4f3cux3092ux4e0eux3048ux308b}}

\[
A_k=\sum_{i=1}^k\sigma_iu_iv_i^T
\]

は rank≤k の中で最良の近似です。

つまり「上位特異値を残す」は便利な経験則ではなく、明確な最適性定理を持っています。

\hypertarget{pca-ux306fux4f55ux3092ux3057ux3066ux3044ux308bux306eux304b}{%
\subsection{PCA
は何をしているのか}\label{pca-ux306fux4f55ux3092ux3057ux3066ux3044ux308bux306eux304b}}

中心化データに対して PCA は

\begin{itemize}
\tightlist
\item
  射影後の分散を最大にする方向を探す
\item
  再構成誤差を最小にする低次元部分空間を探す
\end{itemize}

という二つの等価な問題です。

SVD によってこの二つが一つにつながります。

\hypertarget{ux5727ux7e2eux3068ux30ceux30a4ux30baux9664ux53bb}{%
\subsection{圧縮とノイズ除去}\label{ux5727ux7e2eux3068ux30ceux30a4ux30baux9664ux53bb}}

どちらも小さい特異成分を捨てますが、意味は異なります。

\begin{itemize}
\tightlist
\item
  圧縮: 少数パラメータで近似する
\item
  denoising: 小さい成分を主にノイズと仮定する
\end{itemize}

後者にはモデル仮定があり、弱い信号まで消す危険があります。

\hypertarget{ux5230ux9054ux57faux6e96}{%
\subsection{到達基準}\label{ux5230ux9054ux57faux6e96}}

「次元削減とは単に列数を減らすことではなく、データの主要部分空間を見つけること」と説明できることを目標にします。

\end{document}
